\section{Expanded related work}
\label{app:related-work}
\subsection{Core related works}
\noindent\textbf{Distillation and the OPD loop.}
Knowledge distillation transfers teacher outputs \citep{hinton2015distillingknowledgeneuralnetwork}, including teacher-generated sequences \citep{kim2016sequence}. DAgger, the seminal online imitation algorithm, addresses distribution shift through expert feedback on learner-visited states \citep{ross2011reduction}; \citet{lin2020autoregressiveknowledgedistillationimitation} develop this connection for autoregressive distillation. MiniLLM uses reverse KL and student-generated samples \citep{gu2023minillm},  \citep{agarwal2024onpolicydistillation,piskorz2026onpolicyoffpolicylearningsystematic} study the choice of divergence loss, and \citet{tinker_opd_2025} demonstrate practical OPD for reasoning and personalization. These methods motivate our question: how does dense supervision on student histories produce successful transfer, stagnation, or degeneration? We develop a systematic first-principles account of the coupled loop through which shared-parameter updates change predictions, renew visited histories, and reshape subsequent supervision.

\noindent\textbf{Kernel response laws.}
The neural tangent kernel connects parameter learning to function-space change \citep{jacot2018neural}. \citet{malladi2023kernelbasedviewlanguagemodel} analyze language-model finetuning through kernels, including an Adam extension; \citet{ren2025learning} decompose prediction changes through normalization, empirical-kernel, and loss-gradient terms for SFT and preference finetuning, including offline DPO. \citet{tomihari2026rl} studies update propagation in RL post-training. OPD needs different theoretical and empirical machinery for two reasons. Its on-policy sampling moves the learning distribution itself, so the analysis must track this motion, which is absent from the offline algorithms (SFT and offline DPO) studied by \citet{ren2025learning}. Its supervision is also dense and token-level and can be disconnected from the reasoning outcome (such as accuracy), whereas RL gives only sparse, outcome-level learning signals.

Nevertheless, the fundamental response law of the shared-parameter neural learner is shared with \citet{ren2025learning}. We then follow the induced motion through the rest of the OPD loop: changes at ancestor decisions alter reach probabilities, renew the training source, and thereby change the requests driving future updates. Combining the evolving kernel, routing/drift split, and continuation-value analysis yields a mechanistic account of how optimization changes future supervision and outcomes (\cref{sec:theory,app:derivations}).

\noindent\textbf{Cumulative weight geometry vs the geometry of request conversion.}
\citet{shen2026geometryonpolicydistillation} find low stable rank and early stabilization of dominant subspaces in cumulative weight-matrix changes; constrained-training experiments test those subspaces' functional sufficiency. We instead analyze the current logit Jacobian and the response it delivers to specified teacher requests. The Jacobian can change even when cumulative displacement is low-rank. Frozen histories, coordinates, and requests isolate amplification and reorientation; crossed teacher fields test preferential adaptation in the same kernel (\cref{app:collective-frozen,app:collective-specificity}). Distinguishing local accessibility from collective coherence further explains why individually expressible requests can cancel. These analyses concern how geometry enables or obstructs transfer.

\noindent\textbf{Local fitting vs task performance.}
Policy-gradient theory \citep{sutton2000policy} and cost-to-go-aware imitation learning \citep{ross2014reinforcementimitationlearninginteractive} connect actions to subsequent outcomes. G-OPD interprets distillation through KL-constrained RL and reward extrapolation \citep{yang2026gopd}. \citet{armandpour2026unmaskingonpolicydistillationhelps} use student continuation-success probabilities to derive an ideal local logit direction and measure its alignment with distillation signals (however, they acknowledge that such estimations in long-CoT reasoning can be very noisy). \cite{uemura2026learningoverwritingtheoryselfdistillation} theoretically studies the dynamics of self-distillation for a single-layer neural network, and analyzes how the supervision affects long-horizon traversal of directed acyclic graphs, invoking tools such as future advantage in their analysis. Our analysis connects this type of outcome perspective to \emph{delivered} motion: $\sum_s K_t(x,s)\Gamma_t(s)$ includes the effects of all coupled source prefixes in addition to the request at $x$. Crucially, prefix readouts $\mathbb{E}_{\nu_t}Y_t(x)$ and terminal outcomes $\mathbb{E}_{\Pi_t}Y$ are different objects. At an intermediate cut-position in the autoregressive rollout, terminal change splits into the change in reaching that prefix and the change in the continuations; unrolling the latter and recombining both terms yields the advantage-weighted contribution of every decision (\cref{eq:broadcast-continuations,app:outcome-proof}). This separates reaching a prefix, fitting there, and completing it successfully --the distinction underlying the degeneration mechanism and continuation readiness.

\noindent\textbf{Degeneration involves source feedback as well as supervision quality.}
\citet{fu2026revisitingonpolicydistillationempirical} study unreliable supervision and implementation choices; \citet{jin2026entropyawareonpolicydistillationlanguage} augment high-entropy teacher supervision with forward KL; and \citet{zhu2026manyfaces} analyze signals induced by top-$K$ truncation, renormalization, and stop-gradient conventions. \citet{rethinking_opd_2026} also observe deterioration with prefix depth. Contemporaneous work by \citet{wang2026demystifying} studies teacher mismatch and length exploitation. Our first-principled analysis reveals a phased failure mechanism: early routing makes already-uncertain continuations prevalent, after which substantial fitting can \emph{sharpen those very histories} without restoring useful generation (\cref{app:rd-failure,app:rd-entropy-evidence}). Our routing-drift machinery, coupled with the learning-value decomposition (for terminal-accuracy readout) distinguish the failure mechanism from a naive, uniform entropy-increasing learning process. The important feedback is that the histories made common by earlier updates become the source of later supervision. The prepared learner's different trajectory further shows why source movement itself is not synonymous with failure, and its continuation-readiness helps in exposing useful subsequent learning.

\noindent\textbf{Offline preparation can improve readiness for learning without improving immediate accuracy.}
\citet{rethinking_opd_2026} motivate cold-start SFT by reducing thinking-pattern mismatch and improving overlap. Their \S5.1/Figure~8 studies a strong standalone-SFT regime: the SFT initialization already exceeds the direct-OPD endpoint on all three plotted benchmarks, with little further gains from OPD. We study a different regime, where the student-teacher gap is significant, and preparation leaves starting accuracy near the unprepared student's but changes subsequent OPD from deterioration to substantial improvement. The principled motivation is \emph{continuation readiness}: the outcome analysis and observed failure mechanisms in \cref{sec:broadcast} identify a distinction between increasing a history's reach and executing its continuation successfully. Teacher-generated trajectories provide training on such histories before the student must reach them itself. Our dose study exposes deferred benefits and a sharply nonmonotonic stability pattern at 7B; the prepared-run analysis traces a different subsequent routing/fitting process (\cref{tab:dose-figure,app:rd-prepared}). 

In \emph{contemporaneous work} Simple-OPD \citep{liu2026simpleopd} studies warm-up data, duration, and LoRA parameterization, including benefits from incorrect teacher-compatible traces and differences between SFT and downstream rankings. These findings complement our mechanistic account linking preparation to source feedback and continuation values. The key result is an altered learning trajectory; the preparation dose-response experiments demonstrate this interesting benefit.

\subsection{More related works}
\noindent\textbf{Reasoning interpretability}. Part of the focus of our work is on how LLMs learn to reason (from a teacher model), and our analysis methodology is partly inspired by the discoveries in the reasoning interpretability community. \citet{venhoff2025base} is especially relevant, as they discovered that RL-trained models tend to acquire the ability to manage ``when'' to fire a reasoning mechanism (be it an ``epistemic control'' mechanism such as backtracking or planning, or computation mechanisms like estimating an inequality); \cite{marjanovic2026deepseekr} directly studies the chain-of-thought reasoning mechanisms of DeepSeek-R1 models, which our DeepSeek student models are distilled from, also inspiring our analysis on the epistemic control transmission in OPD. More broadly, \cite{bogdan2025thoughtanchorsllmreasoning, baek2025understandingdistilledreasoningmodels, venhoff2025understandingreasoningthinkinglanguage, troitskii-etal-2025-internal, NEURIPS2025_2c447757, ward2025reasoningfinetuningrepurposeslatentrepresentations} identify causal mechanisms behind LLMs' reasoning behaviors in logic, mathematics and other related domains.

\textbf{OPD and its variants}. On-policy distillation is becoming increasingly popular in industry post-training pipelines \citep{deepseekai2026deepseekv4highlyefficientmilliontoken, kimiteam2026kimik3openfrontier}. Many variants are being proposed, including selecting which tokens to perform supervision on \citep{jin2026entropyaware,yu2026mismatchmattersonpolicydistillation}, and curriculum-styled variants \citep{ziheng2026moreearlystoppingrollout,zhang2026rolloutsnecessaryonpolicydistillation,li2026onpolicydistillationcurriculumturnlevel}. The survey on OPD \citet{song2026surveyonpolicydistillationlarge} is also a good reference for understanding the research landscape in this area.

\textbf{Learning dynamics}. Many classic works lie in this space and are predecessors to recent studies of LLM learning dynamics. \cite{Biehl_1995,PhysRevLett.74.4337,saxe2014exactsolutionsnonlineardynamics,NEURIPS2019_cab070d5} are early representative works, studying how learning trajectories and generalization depend on the model, initialization, and training process, and providing foundations for our analysis of how teacher supervision changes the student over time. Furthermore, a particular line of work study feature learning in nonlinear neural networks: \cite{pmlr-v130-baratin21a,atanasov2022neural,hong2024why,9719702,pmlr-v195-allen-zhu23a} study how the nonlinear neural network’s geometry evolves and aligns with the task, showing feature re-orientation early on in (pre-)training that precedes appreciable loss reduction (sometimes called the ``lottery winning'' phase, where certain neurons or features better align towards loss-reducing directions early in training without substantial loss-reduction yet). These works are especially relevant to our analysis of the student’s changing responsiveness to teacher requests in \cref{sec:conversion}, and to offline preparation whose benefits become apparent through subsequent learning. Additionally, another relevant line of work in the pre-LLM era study dynamics of DNN distillation. \cite{pmlr-v97-phuong19a,NEURIPS2020_ef0d3930,Hong_2021_CVPR,allenzhu2023understandingensembleknowledgedistillation,10.1145/3511808.3557237,he2022feature} examine how the geometry of the learning problem, teacher knowledge, and the student’s representation affect transfer, using linear, deep linear, or kernel-based analyses of the SGD-based learning trajectory.